\section{Semantic Analysis Implementation}\label{sec:semantic-implementation}

This section describes how the semantic rules from Chapter~\ref{ch:language-design} are realized in the
\texttt{motivoc} codebase.
Semantic analysis is compiled into \texttt{motivo\_semantic}.
The public façade \texttt{semantic/analyze.cpp} constructs an \texttt{AnalysisResult} from the parsed program and
runs \texttt{detail::Traversal(result, diagnostics).run(program)}.
Implementation files partition responsibilities as follows.

\noindent \textbf{\texttt{detail/traversal.cpp}.} Defines \texttt{Traversal::run}, header validation
(\texttt{validate\_header}), and \texttt{add\_pattern\_symbol} (registers overload signatures via
\texttt{ScopeStack::add\_symbol}).

\noindent \textbf{\texttt{detail/traversal/collectors.cpp}.} Implements \texttt{collect\_global\_patterns},
\texttt{collect\_track\_patterns}, and \texttt{collect\_voice\_patterns}, each scanning local \texttt{PatternDefinition}
nodes before bodies are visited so calls can resolve overloads.

\noindent \textbf{\texttt{detail/traversal/validators.cpp}.} Hosts \texttt{validate\_binary\_operands},
\texttt{validate\_numeric\_operand}, and \texttt{validate\_call} (signature lookup through
\texttt{find\_pattern\_visible\_by\_signature} plus per-argument \texttt{is\_assignable} checks).

\noindent \textbf{\texttt{detail/traversal/visitors/}.} Split translation units keep compile times manageable:
\texttt{visit\_definitions.cpp} (\texttt{visit\_globals}, \texttt{visit\_track}, \texttt{visit\_voice},
\texttt{visit\_pattern}), \texttt{visit\_statements.cpp} (control flow, declarations, assignments, \texttt{play}),
and \texttt{visit\_expressions.cpp} (literals, operators, sequences, chords, pattern calls).

\noindent \textbf{\texttt{detail/symbol\_table.cpp} and \texttt{detail/scopes/scope\_stack.cpp}.} The symbol table
stores \texttt{Scope} records and \texttt{Symbol} entries indexed by \texttt{SymbolId}; the scope stack pushes/pops
lexical scopes as the visitor enters global, track, voice, pattern, and block nodes.
Pattern overload uses \texttt{find\_in\_current\_scope\_by\_signature} rather than arity alone.

\noindent \textbf{\texttt{detail/annotations.cpp}.} Maps expression pointer addresses to computed \texttt{types::Type}
values and resolved \texttt{SymbolId}s; also records assignment targets for \texttt{AssignStatement} nodes.
Lowering queries these maps through \texttt{AnalysisResult} accessors and never re-runs type inference.

Each \texttt{visit\_*} method returns a \texttt{types::Type} for expressions or \texttt{Void} for statements, writing
results into \texttt{Annotations} before returning so downstream siblings and enclosing expressions observe consistent
types within the same traversal.
